% Draft.
\documentclass[11pt,a4paper]{amsart}

\usepackage[utf8]{inputenc}
\usepackage[T1]{fontenc}
\usepackage{amsmath, amssymb, amsthm}
\usepackage{booktabs}
\usepackage{hyperref}
\hypersetup{hidelinks}
\usepackage{geometry}
\geometry{margin=2.5cm}

\newtheorem{theorem}{Theorem}
\newtheorem{corollary}[theorem]{Corollary}
\theoremstyle{remark}
\newtheorem{remark}[theorem]{Remark}

\newcommand{\Z}{\mathbb{Z}}

\title[A verified Miller test with Bach's constant]{A Lean-verified Miller test with Bach's explicit constant}
\author{Ernesto \'Alvarez}
\thanks{This work is part of the AMate project, by Ernesto \'Alvarez. The formal proofs were developed with the assistance of language models.}
\date{28 November 2026 (draft)}

\begin{document}

\begin{abstract}
We formalize in Lean~4, with Mathlib, Miller's primality test with an explicit Bach-type bound: assuming GRH and several analytic statements about Dirichlet $L$-functions, every non-principal character modulo $n$ takes a value different from $1$ at some $a\le 2(\log n)^2$ coprime to $n$, if $\log n\ge 64$; with constants $16$ and $4$ the threshold drops to $\log n\ge 8.3$ and $\log n\ge 23.5$. From this we obtain an executable Miller test whose correctness is proved for every $n$, a command-line program that computes it, and Bach's classical corollaries: small generators of $(\Z/n\Z)^\times$, small $k$-th power non-residues and a deterministic Solovay--Strassen test. The analytic hypotheses and GRH are theorems of the Lean library announced by Alp\"oge (2026); composing with it removes every hypothesis. Everything depends only on Lean's three standard axioms.
\end{abstract}

\maketitle

\section{Introduction}

Miller~\cite{Miller} proved that, under the generalized Riemann hypothesis (GRH), testing the bases $a\le C(\log n)^2$ decides whether $n$ is prime, and Bach~\cite{Bach} gave the constant $C=2$. Without GRH, the Miller--Rabin test~\cite{Rabin} is probabilistic. What we contribute is not new mathematics but a complete formalization in Lean~4~\cite{Lean4} with Mathlib~\cite{Mathlib}: the algebraic part, the explicit bound, and a test that runs and whose answer is proved. The proof of the constant is our own; it follows the classical scheme of the weighted explicit formula, but it does not reproduce Bach's paper line by line.

Our hypotheses are of two kinds: GRH for primitive characters, and seven analytic statements about $L(s,\chi)$ (Cauchy's formula on rectangles, the location and finiteness of the zeros, the regular part and the functional equation of $L'/L$, bounds for the digamma function, and good heights), plus the representation of the digamma function as a limit. All of them are theorems of the Lean library announced by Alp\"oge~\cite{Alpoge}, which has been machine-checked~\cite{MIT} but, to our knowledge, has no published human review.

\section{The algebraic part}

Let $n\ge3$ be odd. An integer $a$ coprime to $n$ is a \emph{strong liar} if $a^d\equiv1$ or $a^{2^rd}\equiv-1\pmod n$ for some $r<s$, where $n-1=2^sd$ with $d$ odd. If $n$ is prime, every $a$ is one; if $n$ is composite, the strong liars lie in a proper subgroup of $(\Z/n\Z)^\times$. Every proper subgroup lies in the kernel of a non-principal Dirichlet character, by duality for finite abelian groups. Hence, if every non-principal character takes a value different from $1$ at some $a\le B(n)$ coprime to $n$, the test with the bases $a\le B(n)$ is correct.

The same holds for Euler's criterion, the basis of the Solovay--Strassen test~\cite{SS}: $a$ is an \emph{Euler liar} if $a^{(n-1)/2}\equiv\left(\frac an\right)\pmod n$, with the Jacobi symbol. We prove without hypotheses that, for odd composite $n$, the Euler liars form a proper subgroup (\texttt{primo\_de\_euler}): if $p^2\mid n$, a unit of order $p$ does not satisfy $u^{n-1}=1$; otherwise $n=pm$ with $p\nmid m$, and by the Chinese remainder theorem there is an $a$ that is a non-residue modulo $p$ and $\equiv1\pmod m$, so that $\left(\frac an\right)=-1$ while $a^{(n-1)/2}\equiv1\pmod m$.

\section{Bach's bound with an explicit constant}

Let $\chi\ne1$ be a character modulo $n$, and suppose that $\chi(a)=1$ for every $a\le x$ coprime to $n$. Let $\chi^*$ be the associated primitive character, of conductor $q$, $L=\log n$ and $\ell=\log q$. We use the kernel $K(s)=x^s/((s+\alpha)(s+\beta))$ with $\alpha=2/5$, $\beta=9/20$, and the weight $w(m)=((m/x)^\alpha-(m/x)^\beta)/(\beta-\alpha)$ for $m<x$. The weighted explicit formula for $\chi^*$ and for $\zeta$, the Hadamard bound for the sum over the zeros, explicit bounds for the digamma function on the left edge, the residues at $-\alpha$ and $-\beta$, and the contribution of the powers of the primes dividing $n$ give, for every $x\ge1$ (\texttt{desigualdad\_bach}),
\begin{multline*}
\frac{100}{203}\,x\le\frac{\sqrt x}2\Bigl(1.082\,\ell+9.64\Bigr)+\frac{10\ell+672.9}{\sqrt x}\\
+\ell D(x)+24X(x)+51.3+4.01L+13.27(L-\ell),
\end{multline*}
with $D(x)=20(x^{-2/5}-x^{-9/20})$ and $X(x)=20(x^{-2/5}+x^{-9/20})$. With $x=cL^2$ the inequality fails when $L$ is large:

\begin{center}
\begin{tabular}{ccc}
\toprule
constant $c$ & threshold & theorem \\
\midrule
$2$ & $\log n\ge64$ & \texttt{cota\_bach\_explicita}\\
$4$ & $\log n\ge23.5$ ($n\ge2^{34}$) & \texttt{cota\_bach\_4}\\
$16$ & $\log n\ge8.3$ ($n\ge2^{12}$) & \texttt{cota\_bach\_16}\\
\bottomrule
\end{tabular}
\end{center}

\begin{theorem}
Under GRH and the analytic hypotheses, if $\log n\ge64$, every character $\chi\ne1$ modulo $n$ takes a value different from $1$ at some $a\le2(\log n)^2$ coprime to $n$; the same holds with $a\le4(\log n)^2$ if $\log n\ge23.5$ and with $a\le16(\log n)^2$ if $\log n\ge8.3$.
\end{theorem}

\section{Consequences}

\begin{corollary}[Certified test]
Let $B(n)=n$ if $n<4096$; $8(\lfloor\log_2n\rfloor+1)^2$ if $4096\le n<2^{34}$; $2(\lfloor\log_2n\rfloor+1)^2$ if $2^{34}\le n<2^{93}$; and $(\lfloor\log_2n\rfloor+1)^2$ if $n\ge2^{93}$. Under the same hypotheses, Miller's test with the bases $a\le B(n)$ decides the primality of every $n$ (\texttt{esPrimoMiller\_\allowbreak certificado}), and so does Euler's test with those bases (\texttt{test\_euler\_\allowbreak certificado}).
\end{corollary}

\begin{corollary}[Generators]
If $\log n\ge64$, the units represented by integers $a\le2(\log n)^2$ generate $(\Z/n\Z)^\times$ (\texttt{generan\_explicito}); for every $n\ge3$, so do those represented by $a\le B(n)$.
\end{corollary}

\begin{corollary}[Non-residues]
Let $p$ be prime and $k\ge2$ a divisor of $p-1$. If $\log p\ge8.3$, some $a\le16(\log p)^2$ is not a $k$-th power modulo $p$, and some $a\le2(\log p)^2$ is not one if $\log p\ge64$ (\texttt{noPotencia\_16}, \texttt{noPotencia\_explicito}). With $k=2$ this is the bound for the least quadratic non-residue, in the line of Ankeny~\cite{Ankeny}.
\end{corollary}

The three corollaries come from the same observation: every proper subgroup of $(\Z/n\Z)^\times$ misses some small element.

\section{Composition and execution}

The analytic hypotheses are literal copies of statements of Alp\"oge's library, and they are instantiated with its theorems by definitional equality (\texttt{HipotesisAlpoge.lean}, which does not import the main module and compiles in a few seconds). With that library's GRH, all of the above results hold without hypotheses. We built the library from source twice, on two different virtual machines, and both times its $175\,187$ \texttt{.olean} files were byte-identical to the MIT list~\cite{MIT}; the compositions of the Miller test with constant $2$ were compiled on top of the second build. The later compositions (the test for every $n$ and the non-residues) were compiled locally against the library's binaries. All of them depend only on \texttt{propext}, \texttt{Classical.choice} and \texttt{Quot.sound}, and the independent kernel \texttt{nanoda} rechecks the closure of the main theorems with no errors.

The test is executable. The program \texttt{amate-primo} uses a copy of the functions without Mathlib, and \texttt{esPrimo\_eq} proves that it computes exactly the certified test. On a laptop it decides $2^{127}-1$ ($16\,129$ bases) and ten other numbers in $1.6$ seconds in total, and $2^{521}-1$ ($271\,441$ bases) in $111$ seconds.

\section{Limitations}

\begin{itemize}
\item The hypothesis-free version depends on Alp\"oge's library, which is machine-checked but has no published human review.
\item The constants are not optimal: Bach obtains $2(\log n)^2$ for every $n$, and we only for $\log n\ge64$.
\item The correctness of the program is proved for the Lean functions; the native code produced by the compiler is not checked by the kernel.
\item Euler's test has no executable version, because Mathlib's Jacobi symbol is not meant for computation.
\end{itemize}

\begin{thebibliography}{9}
\bibitem{Alpoge} L.~Alp\"oge, untitled announcement note, \url{https://alpo.ge/remark.pdf}, 20 November 2026. The note is unsigned; the attribution comes from the domain and from~\cite{MIT}.
\bibitem{Ankeny} N.~C.~Ankeny, \emph{The least quadratic non residue}, Ann.\ of Math.\ 55 (1952), 65--72.
\bibitem{Bach} E.~Bach, \emph{Explicit bounds for primality testing and related problems}, Math.\ Comp.\ 55 (1990), no.~191, 355--380.
\bibitem{Lean4} L.~de~Moura, S.~Ullrich, \emph{The Lean~4 theorem prover and programming language}, CADE-28, LNCS 12699 (2021), 625--635.
\bibitem{Mathlib} The mathlib Community, \emph{The Lean mathematical library}, CPP 2020, 367--381.
\bibitem{Miller} G.~L.~Miller, \emph{Riemann's hypothesis and tests for primality}, J.\ Comput.\ System Sci.\ 13 (1976), 300--317.
\bibitem{MIT} A.~Sutherland, \emph{Verification files for the Lean proof of the Riemann Hypothesis}, \url{https://math.mit.edu/~drew/grh/}, accessed 23 November 2026.
\bibitem{Rabin} M.~O.~Rabin, \emph{Probabilistic algorithm for testing primality}, J.\ Number Theory 12 (1980), 128--138.
\bibitem{SS} R.~Solovay, V.~Strassen, \emph{A fast Monte-Carlo test for primality}, SIAM J.\ Comput.\ 6 (1977), 84--85.
\end{thebibliography}

\end{document}
